\section*{Chapter \ref{ch:g12:curly-arrows} --- Reaction Mechanisms: Curly Arrows}

\begin{solution}{exo:g12:curly-arrows:1}
\ce{OH-}: donor (\omterm{def:g10:lewis-and-shape:lone-pair}{lone pairs}, negative charge); no acceptor site.
\ce{H+}: acceptor (no \omterm{def:g9:inside-the-atom:nucleus}{electrons} at all). \ce{CH2=CH2}: donor (the \omterm{def:g10:lewis-and-shape:multiple-bond}{double
bond}). \ce{CH3-Cl}: donor (\omterm{def:g10:lewis-and-shape:lone-pair}{lone pairs} of \ce{Cl}), acceptor (the carbon,
$\delta^+$). \ce{H2O}: donor (\omterm{def:g10:lewis-and-shape:lone-pair}{lone pairs} of \ce{O}); its hydrogen \omterm{def:g7:atoms-and-molecules:atom}{atoms},
$\delta^+$, are weak acceptor sites.
\end{solution}

\begin{solution}{exo:g12:curly-arrows:2}
(a) \omterm{def:g12:curly-arrows:categories}{substitution}; (b) \omterm{def:g12:curly-arrows:categories}{addition}; (c) \omterm{def:g12:curly-arrows:categories}{elimination}; (d) \omterm{def:g12:curly-arrows:categories}{addition}.
\end{solution}

\begin{solution}{exo:g12:curly-arrows:3}
It starts from a pair of \omterm{def:g9:inside-the-atom:nucleus}{electrons} (a \omterm{def:g10:lewis-and-shape:lone-pair}{lone pair} or a bond) and ends on
the \omterm{def:g7:atoms-and-molecules:atom}{atom} where the pair goes. From a \omterm{def:g10:lewis-and-shape:lone-pair}{lone pair} of oxygen to a carbon: the
pair becomes a new \ce{O-C} bond.
\end{solution}

\begin{solution}{exo:g12:curly-arrows:4}
Ammonia gives the \omterm{def:g10:lewis-and-shape:lone-pair}{lone pair} of its nitrogen; the arrow goes from that
\omterm{def:g10:lewis-and-shape:lone-pair}{lone pair} to \ce{H+}. The product \ce{NH4+} carries the charge $+1$.
\end{solution}

\begin{solution}{exo:g12:curly-arrows:5}
A species formed in one step and used up in a later one; it is not in
the overall equation. Here the carbocation \ce{CH3-CH2+}.
\end{solution}

\begin{solution}{exo:g12:curly-arrows:6}
One arrow from a \omterm{def:g10:lewis-and-shape:lone-pair}{lone pair} of the oxygen of \ce{OH-} to the carbon; one
arrow from the \ce{C-Br} bond to the bromine. Charges: $-1$ on the left
(\ce{OH-}); on the right, methanol is neutral and \ce{Br-} carries $-1$:
the total is $-1$ on both sides.
\end{solution}

\begin{solution}{exo:g12:curly-arrows:7}
Bromine is more electronegative than hydrogen: it takes the pair and
leaves as \ce{Br-} ($-1$); the \omterm{def:g11:polarity-and-cohesion:hydrogen-bond}{hydrogen bonds} to a carbon, and the other
carbon is left with a positive charge: \ce{CH3-CH2+} ($+1$).
\end{solution}

\begin{solution}{exo:g12:curly-arrows:8}
1.~\ce{CH3-CH2-OH + H+ -> CH3-CH2-OH2+}: donor, a \omterm{def:g10:lewis-and-shape:lone-pair}{lone pair} of the
oxygen; acceptor, \ce{H+}. 2.~\ce{CH3-CH2-OH2+ -> CH3-CH2+ + H2O}: the
\ce{C-O} pair goes to the positive oxygen (acceptor), which leaves with
it as water. 3.~\ce{CH3-CH2+ -> CH2=CH2 + H+}: the pair of a \ce{C-H} bond
(donor) moves to the positive carbon (acceptor) to form the \omterm{def:g10:lewis-and-shape:multiple-bond}{double bond},
and \ce{H+} leaves. The hydrogen \omterm{def:g9:ions:ion}{ion} used in step 1 is given back in step
3: it appears on both sides and cancels.
\end{solution}

\begin{solution}{exo:g12:curly-arrows:9}
An arrow from a \omterm{def:g10:lewis-and-shape:lone-pair}{lone pair} of the oxygen of water to the positive carbon.
After loss of \ce{H+}, ethanol \ce{CH3-CH2-OH}. Overall:
\ce{CH2=CH2 + H2O -> CH3-CH2-OH}, an \omterm{def:g12:curly-arrows:categories}{addition}.
\end{solution}

\begin{solution}{exo:g12:curly-arrows:10}
Two pairs move, in one step: the \ce{C-O} bond forms, the \ce{C-Br}
bond breaks. No intermediate.
\end{solution}

\begin{solution}{exo:g12:curly-arrows:11}
Chlorine is more electronegative ($3.16 - 2.55 = 0.61$ against
$2.96 - 2.55 = 0.41$): the carbon of chloromethane is the more
$\delta^+$, which alone would make it react faster. It does not: the
speed also depends on how easily the \omterm{def:g10:electron-shells:family}{halogen} leaves with the pair, and
the larger bromine leaves more easily.
\end{solution}

\begin{solution}{exo:g12:curly-arrows:12}
If the \omterm{def:g11:polarity-and-cohesion:hydrogen-bond}{hydrogen bonds} to carbon 1: \ce{CH3-CH+-CH3}, which gives
2-chloropropane \ce{CH3-CHCl-CH3}. If it bonds to carbon 2:
\ce{+CH2-CH2-CH3}, which gives 1-chloropropane \ce{CH2Cl-CH2-CH3}.
\end{solution}

\begin{solution}{exo:g12:curly-arrows:13}
The arrow goes the wrong way: it must start from the donor, a \omterm{def:g10:lewis-and-shape:lone-pair}{lone pair}
of \ce{OH-}, and end on the acceptor carbon. And the new bond is between
oxygen and carbon, not oxygen and bromine: bromine leaves, taking the
pair of the \ce{C-Br} bond.
\end{solution}

\begin{solution}{exo:g12:curly-arrows:14}
The carbon of the \ce{C=O} group of an acid is bonded to two
electronegative oxygen \omterm{def:g7:atoms-and-molecules:atom}{atoms}: strongly $\delta^+$, an acceptor. The
oxygen of the \omterm{def:g11:functional-groups:alcohol}{alcohol} carries $\delta^-$ and two \omterm{def:g10:lewis-and-shape:lone-pair}{lone pairs}: a donor.
Water is lost overall; the \ce{-OH} of the acid is replaced by the
\ce{-O-R} of the \omterm{def:g11:functional-groups:alcohol}{alcohol}: a \omterm{def:g12:curly-arrows:categories}{substitution}.
\end{solution}

\begin{solution}{exo:g12:curly-arrows:15}
Pushed away by the \omterm{def:g9:inside-the-atom:nucleus}{electrons} of the \omterm{def:g10:lewis-and-shape:multiple-bond}{double bond}, the pair of the
\ce{Br-Br} bond moves towards the far bromine: the near bromine becomes
$\delta^+$, an acceptor site. First step: an arrow from the \omterm{def:g10:lewis-and-shape:multiple-bond}{double bond}
to the near bromine, and one from the \ce{Br-Br} bond to the far
bromine, which leaves as \ce{Br-}; a positive carbon is left on the
other carbon of the former \omterm{def:g10:lewis-and-shape:multiple-bond}{double bond}.
\end{solution}

\begin{solution}{pb:g12:curly-arrows:1}
\textbf{1.} Ethene: \ce{H2C=CH2}, a \omterm{def:g10:lewis-and-shape:multiple-bond}{double bond} and four \ce{C-H}.
Hydrogen bromide: \ce{H-Br} with three \omterm{def:g10:lewis-and-shape:lone-pair}{lone pairs} on bromine.

\textbf{2.} The \omterm{def:g10:lewis-and-shape:multiple-bond}{double bond}: its second pair is a region rich in
\omterm{def:g9:inside-the-atom:nucleus}{electrons}.

\textbf{3.} $2.96 - 2.20 = 0.76$: \ce{H} carries $\delta^+$.

\textbf{4.} The hydrogen \omterm{def:g7:atoms-and-molecules:atom}{atom}.

\textbf{5.} An arrow from the \omterm{def:g10:lewis-and-shape:multiple-bond}{double bond} to the hydrogen of \ce{HBr};
an arrow from the \ce{H-Br} bond to the bromine.

\textbf{6.} The carbocation \ce{CH3-CH2+} ($+1$) and \ce{Br-} ($-1$).

\textbf{7.} Before: $0 + 0 = 0$; after: $+1 - 1 = 0$.

\textbf{8.} An arrow from a \omterm{def:g10:lewis-and-shape:lone-pair}{lone pair} of \ce{Br-} to the positive
carbon.

\textbf{9.} Six (three bonds): one pair short of an octet, so it accepts
a pair from any donor nearby.

\textbf{10.} The carbocation \ce{CH3-CH2+}.

\textbf{11.} It is formed in the first step and used up in the second.

\textbf{12.} An \omterm{def:g12:curly-arrows:categories}{addition}.

\textbf{13.} Only the group: the chain of two carbons is kept, and the
\omterm{def:g10:lewis-and-shape:multiple-bond}{double bond} becomes a single bond carrying \ce{H} and \ce{Br}.

\textbf{14.} \qty{28.0}{g/mol}; \qty{108.9}{g/mol}.

\textbf{15.} $10.0 / 28.0 = \qty{0.357}{mol}$.

\textbf{16.} Ethene: $0.357 - x$; \ce{HBr}: in excess; bromoethane: $x$.
Ethene is limiting: $x_{\max} = \qty{0.357}{mol}$.

\textbf{17.} $0.357 \times 108.9 = \qty{38.9}{g}$.

\textbf{18.} $0.75 \times 38.9 = \qty{29.2}{g}$.

\textbf{19.} About \qty{29}{g} of bromoethane.
\end{solution}
